\section{Reflection positivity and the plane covariance}
\label{sec:spectral}

This section identifies every subsequential plane covariance, without assuming
either Gaussianity or rotation invariance of a scaling limit. The positive
transfer method of \cite[Section~5]{Triangular} suggests the argument, but its
sixfold rotational cancellation is unavailable here. Two square-lattice
reflection frames at angle $\pi/4$ provide the required replacement.

We temporarily choose Euclidean units in which the corner grid is $\Z^2$.
This is a fixed rotation and dilation of the original embedding. The local
plaquette calculation above identifies its height weights with the
six-vertex weights $(1,1,c)$. The spin and decoration law obtained from flat
exhaustions is the measure of \cite[Theorem~17.1]{DKLM2026};
\cite[Lemma~18.7(iii)]{DKLM2026} identifies its height-gradient marginal
with the slope-zero six-vertex plane measure. By
\cite[Theorem~2.2]{DKLM2026}, that marginal is stationary under all
corner-grid translations and invariant under the lattice reflections.
The geometric height reflection used below sends $K$ to its composition
with the spatial reflection, up to an additive constant. In arrow variables
this combines spatial reflection with global arrow reversal; it also
preserves the finite plaquette weights directly.
These statements apply for every $1\le c\le2$. A uniform simultaneous
choice of the two spin signs is understood when residues rather than
increments are observed. In particular an odd corner-grid translation $T$
exchanges the parities with gauge
$(\tau',\tau^{\dagger\prime})=(\tau^\dagger\circ T,-\tau\circ T)$:
its height is $K'(x)=K(Tx)-K(T0)$. This preserves the incident-corner rule
and accounts for the change of reference parity.
The additive constant of $K$ is immaterial. For a smooth compactly supported
function $f$ of integral zero, use bilinear nodal interpolation on each
corner-grid square to define $K_\delta(f)$. This interpolation commutes with
all four row reflections and with grid translations.
Write its finite array of coefficients as $b_{\delta,x}$; their sum is
exactly zero. For a $2m$th moment, choose smooth mass-one reference
densities in $2m$ disjoint fixed disks outside the common support of the
tests. Let $r_{\delta,j,y}$ be nonnegative discretizations of these
densities, normalized so that $\sum_y r_{\delta,j,y}=1$. For each occurrence
$j$ the exact finite-sum identity is
\[
 K_\delta(f)
   =\sum_{x,y}b_{\delta,x}r_{\delta,j,y}\bigl(K(x)-K(y)\bigr).
\]
Use a separate reference index $j$ in each factor of the moment. The
regularity estimate \cite[Theorem~4.5]{DKLM2026} then applies to increment
correlations with separately summed reference variables. Its collision
singularities are products of logarithms, which are locally integrable;
the lattice cells of exact coincidence contribute at most a lattice area
times a fixed power of $1+|\log\delta|$. Integrating the estimate gives
\begin{equation}
 \sup_{0<\delta<\delta_0}\E|K_\delta(f)|^{2m}<\infty,
 \qquad m\in\N.
 \label{eq:plane-neutral-moments}
\end{equation}
The bounds are uniform for tests in a bounded smooth family with common
compact support. The estimate holds throughout $1\le c\le2$, not just the
range of the convergence theorem in that paper.
The elementary path-factor correction in its proof is supplied in
\cref{sec:input-path}.

\subsection{Four positive reflection frames}

\begin{lemma}[Positive normal transfer]
\label{lem:positive-transfer}
The plane law is reflection positive across a row of corner-grid sites in
each of the directions $0,\pi/4,\pi/2,3\pi/4$. In its reflection Hilbert
space, translation by two successive row spacings in the inward normal
direction is a positive self-adjoint contraction. For the height-increment
law, tangential translation by one row period is a commuting unitary.
Reflection positivity and positive normal transfer remain valid for local
spin and edge decorations strictly on one side of the row, and for
independent uniform phase lifts of the height. With a fixed one-parity edge
decoration, the tangential unitary assertion is restricted to
parity-preserving translations.
\end{lemma}

\begin{proof}
First use a finite reflection-symmetric flat exhaustion. Across an axial row,
plaquette interactions on the two sides factor after conditioning on the row
heights. Across a diagonal row, a straddling plaquette has the following
conditional weight matrix. If its two row heights are equal, the other
heights are independently one unit above or below that value. Equal
transverse heights have weight $c$, and unequal ones have weight one.
Consequently the matrix, after multiplication by $t$, is
\[
 \begin{pmatrix}1&t\\t&1\end{pmatrix}.
\]
It is positive semidefinite because $0<t\le1$. If the row heights differ
by two, the transverse heights are forced and there is only a scalar factor.
Factoring each positive matrix as a Gram matrix proves positivity after
summing the row data. Independent local decorations are included in the
factors on their own sides. Spin residues are determined by heights and one
compatible global sign choice, which can be fixed at an anchor on the row.
For a phase $U\exp(i\alpha(K(x)-K(o)))$, also condition on its value at that
anchor; $U$ has independent uniform law. Thus these augmentations preserve
the same Gram representation. Circuit comparison passes each reflected
exhaustion to the same plane law.

Let $\Theta$ be reflection in the row and set
$\langle F,G\rangle=\E[\overline{\Theta F}G]$ for functions on its inward
half-plane. Denote the pure normal translation through two row spacings by
$T$. Stationarity gives symmetry of $T$ for this form. Reflection positivity
at the intermediate row gives $\langle F,TF\rangle\ge0$. To prove boundedness
before quotienting, put $b_j=\|T^jF\|$. Reflection Cauchy--Schwarz gives
$b_j^2\le b_{j-1}b_{j+1}$, whereas ordinary Cauchy--Schwarz and stationarity
give $b_j\le\|F\|_{L^2}$ for all $j$. A bounded nonnegative log-convex
sequence cannot increase at its first step. Hence $\|TF\|\le\|F\|$, also
when $\|F\|=0$. Thus $T$ descends to a positive self-adjoint contraction.
Tangential translation preserves the form, has a form-preserving inverse,
and commutes with $T$, whenever it preserves the chosen state space. For
the height-increment law this includes one row period. With a fixed
one-parity decoration it includes two axial periods or one diagonal period;
the normal transfers used above already preserve that parity.
\end{proof}

In a frame with unit tangent $u$ and normal $v$, write the rectangular
sublattice just obtained as
\[
 \Lambda_0=\ell\Z u+d\Z v.
\]
For an axial frame one may take $(\ell,d)=(1,2)$; for a diagonal frame take
$(\ell,d)=(\sqrt2,\sqrt2)$. Each has finite index in the corner grid. All
frequency coordinates below are Euclidean: $p=ku+lv$. In particular, the
normal spectral scale and $|k|$ have the same physical units.
The spectral argument uses the undecorated height-increment space; the
decorated argument in \cref{sec:screening} uses only normal transfers.

\subsection{The spectral measure and its Cauchy disintegration}

\begin{lemma}[Increment spectral measure]
\label{lem:increment-spectral}
There is a nonnegative measure $C$, locally finite on the punctured reciprocal
torus of the corner grid, such that every finitely supported real array of
zero sum satisfies
\begin{equation}
 \E\left(\sum_x b_xK(x)\right)^2
   =\int\left|\sum_x b_x e^{ip\cdot x}\right|^2 C(\dd p).
 \label{eq:increment-spectral}
\end{equation}
There is no additional random linear tilt. For sufficiently small $r$,
\begin{equation}
 C\{r\le|p|\le2r\}\le C_0.
 \label{eq:spectral-annulus}
\end{equation}
Consequently $\int_{0<|p|<r_0}|p|^\eta C(\dd p)<\infty$ for every $\eta>0$.
\end{lemma}

\begin{proof}
Let $V_x$ be ordinary translation on the plane probability space, $E(\dd p)$
its spectral resolution \cite[Theorem~5.6]{WilliamsSpectral2018},
and $\xi_x=K(x)-K(0)$. The cocycle identity gives
\[
 (V_y-I)\xi_x=(V_x-I)\xi_y.
\]
On any spectral neighborhood where $e^{ip\cdot y}-1\ne0$, define
\[
 C(\dd p)=
 \frac{\langle\xi_y,E(\dd p)\xi_y\rangle_{L^2}}
 {|e^{ip\cdot y}-1|^2}.
\]
The identity makes the definitions agree on overlaps; two grid directions
cover the punctured torus. The invariant part is additive and has the form
$A\cdot x$ for a random vector $A\in L^2$. Smearing at scale $L$ against a
neutral test of first moment $m$ gives invariant part $LA\cdot m$.
Its norm is uniformly bounded by \cref{eq:plane-neutral-moments}, since
orthogonal projection cannot increase norm. Two independent choices of $m$
give $A=0$.

Choose two real neutral smooth tests whose Fourier transforms have no common
zero on $1\le|p|\le2$; derivatives in the two coordinate directions of a
sufficiently narrow smooth bump suffice. Their scale-$r^{-1}$ interpolation
transforms converge uniformly on this annulus to those transforms. The sum
of their squared absolute values is therefore bounded below on
$r\le|p|\le2r$. Apply \cref{eq:increment-spectral} and the uniform second
moment bound. Summing over dyadic annuli proves the last assertion.
\end{proof}

Project $C$ to the reciprocal rectangle
$[-\pi/\ell,\pi/\ell)\times[-\pi/d,\pi/d)$ of $\Lambda_0$, and call the
result $C'$. Besides the physical origin, the preimages of its origin are
finitely many nonzero reciprocal points, called aliases. Their neighborhoods
have finite mass; atoms at aliases are allowed.
Put $P=2\pi/d$. For $0<s<\infty$, define the probability law
\[
 Q_s(\dd l)=\sum_{j\in\Z}
 \frac{s\,\dd l}{\pi(s^2+(l+jP)^2)},\qquad -P/2\le l<P/2,
\]
with $Q_0=\delta_0$ and $Q_\infty$ uniform on the circle.

\begin{lemma}[Positive Cauchy mixture]
\label{lem:cauchy-mixture}
For $k\ne0$ there is a nonnegative mixing measure $\nu$ such that
\begin{equation}
 C'(\dd k,\dd l)=\int_{[0,\infty]}Q_s(\dd l)\,\nu(\dd k,\dd s).
 \label{eq:cauchy-mixture}
\end{equation}
It is finite on horizontal bands away from zero and satisfies
\begin{equation}
 \int_{0<|k|<\kappa}|k|\,\nu(\dd k,\dd s)<\infty
 \label{eq:weighted-mixture}
\end{equation}
for small fixed $\kappa$.
\end{lemma}

\begin{proof}
The row increment $F=K(\ell u)-K(0)$ belongs to the reflection Hilbert space.
The joint spectral theorem for its commuting tangential unitary and positive
normal contraction \cite[Theorem~5.6]{WilliamsSpectral2018} gives,
for $j\ge0$,
\[
 \E[F V_{n\ell u+jdv}F]
   =\int e^{in\ell k}\lambda^j\,\gamma(\dd k,\dd\lambda),
 \qquad 0\le\lambda\le1.
\]
Reflection gives the same expression with $|j|$. Set $\lambda=e^{-ds}$,
including $s=\infty$ when $\lambda=0$. These are precisely the vertical
Fourier coefficients of $Q_s$. Uniqueness of Fourier coefficients for finite
measures identifies the ordinary increment measure as this mixture. That
measure is $|e^{i\ell k}-1|^2C'$. Division for $k\ne0$ proves the formula.
Since every $Q_s$ has mass one, the integral in
\cref{eq:weighted-mixture} equals $\int|k|C'(\dd k,\dd l)$ over the same
band. Near the physical origin it is finite by
\cref{eq:spectral-annulus}; all other contributions have finite mass.
\end{proof}

\subsection{An angular defect in two frames}

Define
\[
 \chi(k,l)=\cos(4\arg(k+il))
       =1-\frac{8k^2l^2}{(k^2+l^2)^2},\qquad
 D(k,s)=\left(\frac{|k|-s}{|k|+s}\right)^2,
\]
where $D(k,\infty)=1$. Uniformly in $s\in[0,\infty]$,
\begin{equation}
 \int\chi(k,l)Q_s(\dd l)=D(k,s)+O(|k|),\qquad k\longrightarrow0.
 \label{eq:angular-poisson}
\end{equation}
For completeness, with $b=|k|$, the bounded analytic function
$((b+iz)/(b-iz))^2$ on the upper half-plane has boundary real part
$\chi(b,l)$. Its Poisson average at $is$ is $((b-s)/(b+s))^2$.
To periodize, use
\[
 1-\chi(b,l)=\frac{8b^2l^2}{(b^2+l^2)^2},\qquad
 \int_\R(1-\chi(b,l))\,\dd l=4\pi b.
\]
The sum of noncentral Cauchy densities over one period is uniformly bounded
in $s$. Indeed split its sum at $|j|$ comparable to $s/P$ and use $1/s$
for the nearer terms and a summable square tail for the rest. The omitted
unfolded tail of $1-\chi$ is $O(b^2)$. This proves
\cref{eq:angular-poisson}, including its two limiting laws.

\begin{proposition}[Summable dispersion defect]
\label{prop:dispersion-defect}
In each of the four reflection frames,
\begin{equation}
 \int_{0<|k|<\kappa}D(k,s)\,\nu(\dd k,\dd s)<\infty,
 \qquad
 \nu\{r\le|k|\le2r\}\le C_1.
 \label{eq:dispersion-defect}
\end{equation}
The second bound holds for every $r>0$.
\end{proposition}

\begin{proof}
Integrate \cref{eq:angular-poisson} over $\eps<|k|<\kappa$.
Its error is bounded independently of $\eps$ by
\cref{eq:weighted-mixture}. Replacing $C'$ outside the small physical disk
by the original measure costs bounded outer and alias mass. In the radial
annulus $\eps<|p|<\kappa$, the omitted endcap $|k|\le\eps$ has negative
fourth harmonic only where $|p|\le C\eps$: negativity implies
$|l/k|<1+\sqrt2$. Its negative part is bounded by a fixed number of the
annuli in \cref{eq:spectral-annulus}. Thus, in a frame of angle $\theta$,
\begin{equation}
 \int_{\eps<|k|<\kappa}D\,\dd\nu
 \le \int_{\eps<|p|<\kappa}
        \cos(4(\arg p-\theta))\,C(\dd p)+C_2.
 \label{eq:frame-defect-bound}
\end{equation}
Choose the same physical inner and outer radii for frames $\theta$ and
$\theta+\pi/4$, making the fixed outer radius smaller if necessary.
Their radial integrands are negatives pointwise. Adding the inequalities
therefore bounds the sum of the two nonnegative defect integrals. Monotone
convergence proves finiteness for each. Apply the same argument to the other
frame pair. No rotation of the measure $C$ through $\pi/4$ has been used.

On $r\le|k|\le2r$, the part with $s\notin[r/2,4r]$ has $D\ge1/9$, so
its mass is bounded. For $s\in[r/2,4r]$, a fixed fraction of $Q_s$ lies in
$|l|\le2r$. The remaining mixing mass is bounded by a constant times
$C'\{r\le|k|\le2r,\ |l|\le2r\}$, which is bounded by the annular estimate
and finite alias mass. This proves the horizontal-band bound for sufficiently
small $r$. For the remaining $r$, the bands lie in a fixed horizontal region
bounded away from zero, whose $\nu$ mass is finite by
\cref{lem:cauchy-mixture}; bands beyond the reciprocal interval are empty.
Increasing $C_1$ gives the bound for every $r>0$.
\end{proof}

\begin{lemma}[Spectral endcaps]
\label{lem:spectral-endcap}
For every frame, sufficiently small $r$, and $0<\eps<r/10$,
\[
 C\{r<|p|<2r,\ |k|\le\eps\}\le C_3\eps/r.
\]
In particular the punctured small disk has no mass on a frame's line $k=0$.
\end{lemma}

\begin{proof}
Describe the same region in a frame at angle $\pi/4$. Its new coordinates
$k',l'$ both have size comparable to $r$. At fixed $k'$, the old restriction
$|k|\le\eps$ confines $l'$ to an interval of length $O(\eps)$. On this
interval each central Cauchy density is at most $C/r$, since
$s/(s^2+l'^2)\le1/(2|l'|)$, and the folded tails are bounded.
The $s=0$ atom is outside the interval; the uniform law also satisfies the
bound. The horizontal-band mass in the new frame is bounded by
\cref{eq:dispersion-defect}. Projection only adds nonnegative alias mass,
so its bound bounds the original physical measure. Letting $\eps$ decrease
to zero gives the last statement.
\end{proof}

\subsection{Green covariance and reflection nullity}

For neutral smooth tests write
\[
 \mathcal L(f,g)=\iint f(x)\log\frac1{|x-y|}g(y)\,\dd x\,\dd y.
\]

\begin{proposition}[Subsequential plane Green covariance]
\label{prop:plane-green}
Every mesh sequence has a subsequence and a finite positive number $\beta$
such that, for all neutral smooth compactly supported tests,
\begin{equation}
 \E[K_\delta(f)K_\delta(g)]\longrightarrow\beta\mathcal L(f,g).
 \label{eq:plane-green}
\end{equation}
The possible values of $\beta$ are bounded above and bounded away from zero.
The convergence also holds for smoothly converging tests in a common compact
set. At this stage $\beta$ is allowed to depend on the subsequence.
\end{proposition}

\begin{proof}
Rescale the physical measure by $p\mapsto p/\delta$. The annular bound gives
vague subsequential compactness off zero; write $M$ for one such limit.
Fix $0<b<B$ in one reflection frame. Let $\nu_\delta^{b,B}$ be the
pushforward of $\ind_{\{b\delta\le |k|\le B\delta\}}\nu$ under
$(k,s)\mapsto(\xi,\sigma)=(k/\delta,s/\delta)$, with $\infty/\delta=\infty$.
The horizontal-band bound in \cref{eq:dispersion-defect} bounds its total
mass uniformly. Homogeneity of $D$ and its integrability there give
\[
 \int D(\xi,\sigma)\,\nu_\delta^{b,B}(\dd\xi,\dd\sigma)
   =\int_{b\delta\le |k|\le B\delta}D(k,s)\,\nu(\dd k,\dd s)
   \longrightarrow0.
\]
Take a weakly convergent further subsequence on the compact space
$\{b\le|\xi|\le B\}\times[0,\infty]$, where the last interval is
compactified at infinity. On this space $D$ is continuous, equals one
at $\sigma=0$ and $\sigma=\infty$, and vanishes exactly at
$\sigma=|\xi|$. Every limiting mixing measure is therefore supported
on that graph. In particular, the defect estimate excludes mass at
both ends of the rescaled normal parameter.

Let $\widehat Q_{\delta,\sigma}$ be the pushforward of $Q_{\delta\sigma}$
under $l\mapsto l/\delta$. On a compact vertical interval, uniformly for
$\sigma$ in a compact subset of $(0,\infty)$,
\[
 \widehat Q_{\delta,\sigma}(\dd l)
   =\left(\frac{\sigma}{\pi(\sigma^2+l^2)}+O(\delta)\right)\dd l.
\]
The first term is the rescaled central Cauchy density; the bound on the
folded tails used in \cref{eq:angular-poisson} gives the error. The vanishing
defect makes the mixing mass outside any fixed neighborhood of
$\sigma=|\xi|$ tend to zero. We may therefore pass this kernel formula
against compactly supported vertical tests and tangential tests in the
interior of the closed band. This restriction avoids any issue from atoms
on its boundary.

For these tests the projected measure $C'$ and the physical measure have
the same limit. Indeed, the preimage near an alias $p_a$ lies in
$B(p_a,c_*\delta)\setminus\{p_a\}$ for a constant $c_*$ depending on the
test support: the nonzero scaled tangential coordinate excludes $p_a$
itself. The measure $C$ is finite near $p_a$, so its mass on this shrinking
punctured neighborhood tends to zero. There are only finitely many aliases.
Exhausting $k\ne0$ by interiors of the bands above now gives
\[
 M(\dd k,\dd l)=\frac{|k|}{\pi(k^2+l^2)}\,\eta(\dd k)\,\dd l
\]
for a nonnegative measure $\eta$. The tangential measures agree on
overlapping bands, since integrating the displayed positive density over
a fixed vertical interval recovers them uniquely.

Thus $|p|^2M$ is locally invariant under normal translations where
$k\ne0$, without assuming that $\eta$ has a density.
At every nonzero point at least two of the four frames have nonzero
tangential coordinate and independent normals. The two corresponding
distributional derivatives of $|p|^2M$ vanish in a neighborhood. Convolving
locally with a smooth kernel shows it is a constant multiple of Lebesgue
measure there. Constants agree on overlaps in the connected punctured plane.
We obtain $M(\dd p)=b_0|p|^{-2}\dd p$, with $0\le b_0<\infty$.

To pass from vague convergence to tests, let $A_{\delta,f}$ be the transform
of the interpolation coefficients. They have exactly zero sum. Near zero,
for every fixed integer $q\ge3$,
\[
 |A_{\delta,f}(p)|\le C_q
 \min\{|p|/\delta,(1+|p|/\delta)^{-q}\}.
\]
The first bound uses the first absolute moment of the coefficients; the
second follows by $q$ lattice summations by parts against the smooth test
convolved with its nodal hat. Away from zero the same summations give
$O(\delta^q)$. Thus the low-frequency tail is $O(\eps^2)$, the high-frequency
tail within the small disk is $O(R^{-2q})$, and the remaining tail vanishes
by local finiteness. On the intervening compact annulus use vague convergence.
The Fourier convention $\widehat f(p)=\int e^{-ip\cdot x}f(x)\dd x$ gives
\[
 b_0\int\frac{\widehat f(p)\overline{\widehat g(p)}}{|p|^2}\dd p
   =2\pi b_0\mathcal L(f,g).
\]
This proves the assertion with $\beta=2\pi b_0$, including varying tests.

Finally choose two disjoint small disks and a difference of nonnegative
mass-one tests in them. Choose two fixed nested annuli enclosing the first
disk and avoiding the second, with positive margins. For all sufficiently
small meshes, the nodal supports remain in their respective disks.
Apply \cref{prop:circuits} successively in the two annuli. With probability
bounded below, they contain $W$ circuits of opposite prescribed colors.
This uses the conditional one-parity bound, so it does not require positive
association between the two opposite-color events.

Reveal the outer circuit by an outside-in stopping search and condition on
it and its exterior. No $N$ cluster crosses this cut. Conditional also on
the complete $(s,W)$ data in its finite interior, \cref{lem:finite-cuts}
gives independent fair colors to the enclosed $N$ clusters. The opposite
$s$ colors of the two circuits imply that an odd number of their intervening
cluster jumps are nonzero. At least one such free cluster has an odd hole
containing the first support, while the second support is outside the
outer cut. Flipping its fair color shifts the first average by two in one
direction or the other relative to the second, and hence changes their
difference by four. Its conditional variance is four.
The resulting strictly positive lower variance bound is uniform in mesh.
It bounds every possible $\beta$ away from zero. The upper bound follows
from \cref{eq:plane-neutral-moments} applied to these tests.
\end{proof}

\begin{corollary}[Reflection-null Laplace insertion]
\label{cor:reflection-null}
Let $\phi\in C_c^\infty(\R^2)$ be supported a positive distance to one side
of a line in any of the four reflection frames. For site rows converging to
that line, the reflection norm of $X_\delta=K_\delta(\Delta\phi)$ tends to
zero. The assertion includes converging placements and the local
augmentations in \cref{lem:positive-transfer}.
\end{corollary}

\begin{proof}
Every subsequence has a further covariance limit in \cref{prop:plane-green}.
The squared reflection norm then tends to the Green pairing of
$\Delta\phi$ with its reflected copy, times $\beta$. This pairing is zero:
after integration by parts it is a constant times the integral of the two
gradients, whose supports are disjoint. The coefficient need not yet be
unique. Positivity and the subsequence argument prove convergence to zero
along the full mesh limit. Decorations not used by $X_\delta$ do not change
this covariance, and converging placements are covered by the varying-test
part of \cref{prop:plane-green}.
\end{proof}
